\documentclass[aspectratio=169,10pt,fontset=fandol]{ctexbeamer}
\usepackage{amsmath,amssymb,mathtools,bm,mathrsfs}
\usepackage{graphicx,booktabs,array,tabularx,multirow}
\usepackage{tikz,pgfplots,xparse,xstring,mdframed}
\usetikzlibrary{arrows.meta,calc,angles,positioning,shapes.geometric,graphs}
\pgfplotsset{compat=1.18}

\setbeamerfont{normal text}{size=\normalsize}
\setbeamerfont{title}{size=\large}
\setbeamerfont{subtitle}{size=\normalsize}
\setbeamerfont{frametitle}{size=\large}
\setbeamerfont{framesubtitle}{size=\small}
\setbeamerfont{block title}{size=\normalsize}
\setbeamerfont{block body}{size=\normalsize}
\setbeamerfont{section in toc}{size=\normalsize}
\setbeamerfont{section in head/foot}{size=\footnotesize}
\setbeamerfont{author in head/foot}{size=\footnotesize}
\setbeamerfont{title in head/foot}{size=\footnotesize}
\setbeamerfont{date in head/foot}{size=\footnotesize}
\setbeamertemplate{navigation symbols}{}
\setbeamertemplate{itemize item}{\normalsize\textbullet}
\setbeamertemplate{itemize subitem}{\small\textbullet}
\setbeamercolor{structure}{fg=blue!65!black}
\setbeamercolor{block title}{fg=blue!65!black,bg=blue!8}
\setbeamercolor{block body}{fg=black,bg=blue!3}
\setbeamertemplate{footline}{\leavevmode\hbox to\paperwidth{\hfill{\footnotesize\insertframenumber}\hspace*{1em}}\vskip3pt}

% ctex fontset=fandol; system fonts are not redistributed.

\newenvironment{proofdetail}{\par\noindent\textbf{证明：}\quad}{\par}
\renewenvironment{solution}{\par\noindent\textbf{解：}\quad}{\par}
\title{书目回归.tex}
\begin{document}
\section{引用来源核对}
% FRAME-ID F001
\begin{frame}[t,fragile]{目录}
\normalsize
{\color{blue}\normalsize 1. 引用来源核对\par}\vspace{0.22em}
\end{frame}
% FRAME-ID F002
\begin{frame}[t,fragile]{引用来源核对}
\normalsize
\begin{block}{定义}
{\normalsize
% SOURCE-BEGIN U0001 0 21
函数是两个集合之间满足唯一对应条件的映射。
% SOURCE-END U0001
}
\end{block}
\end{frame}
% FRAME-ID F003
\begin{frame}[t,fragile]{引用来源核对}
\normalsize
{\normalsize
% SOURCE-BEGIN U0002 0 35

土壤性质的背景材料见文献\cite{BradyWeil}。


% SOURCE-END U0002
}
\end{frame}
% FRAME-ID F004
\begin{frame}[t,fragile]{参考文献}
\normalsize
{\normalsize
% SOURCE-BEGIN U0003 0 172
\begin{thebibliography}{99}
\bibitem{BradyWeil} WEIL R R, BRADY N C. \emph{The Nature and Properties of Soils}（15th ed.）[M]. Harlow: Pearson, 2017.

\end{thebibliography}
% SOURCE-END U0003
}
\end{frame}
\end{document}
